\documentclass[10pt,a4paper]{amsart}
\usepackage[margin=19mm]{geometry}
\usepackage{amsmath,amssymb,mathtools}
\usepackage[scr=rsfs,cal=euler]{mathalfa}
\usepackage{xfrac}
\usepackage[unicode,hidelinks,bookmarks=false]{hyperref}
\newcommand{\ie}{i.e.\ }
\newcommand{\cf}{cf.\ }
\newcommand{\resp}{respectively\ }
\newcommand{\Subsection}{\subsection}
\providecommand{\nobreakdash}{\nobreak\hbox{-}}
\setlength{\emergencystretch}{2em}
\multlinegap0pt
\usepackage{amssymb}
\usepackage{enumerate}
\usepackage[all]{xy}
\newtheorem{definition}{Definition}
\newtheorem{proposition}{Proposition}
\newcommand{\F}{{\mathbb{F}}}
\newcommand{\N}{{\mathbb{N}}}
\newcommand{\SSigma}{{\overline{\Sigma}}}
\newcommand{\VI}{{\mathrm{VI}}}
\title{Bounding regularity of $\VI^m$-modules}
\author{Wee Liang Gan, Khoa Ta}
\date{Source: arXiv:2512.25010v1 (CC BY 4.0)}
\begin{document}
\maketitle

\noindent\emph{Source and licence.} Bounding regularity of $\VI^m$-modules. Version arXiv:2512.25010v1 (CC BY 4.0), distributed by arXiv under the Creative Commons Attribution 4.0 International licence, http://creativecommons.org/licenses/by/4.0/, as recorded in the arXiv licence metadata for this exact version and shown on the abstract page https://arxiv.org/abs/2512.25010v1. Full licence notice: \texttt{licenses/arxiv-2512.25010-CC-BY-4.0.txt}.\par\smallskip

\noindent\textit{Editorial scope.} Counted unit: the article's proposition prop:barshift1 (source0.tex lines 413--421). The article's complete original proof of it is delivered with it (source0.tex lines 423--442, one \texttt{proof} environment of 20 lines). No mathematics is written, repaired or re-derived: the statement and the proof are the article's own lines, in the article's own notation, and no retained line is substituted or re-flowed.  Retained uncounted as the source-local context the proof needs: lines 405--411.  Nothing was omitted from any retained span, so the package carries no omission record.  No retained line contains a citation, so the wrapper carries no bibliography and no bibliography entry.  No retained line contains a figure environment, an image-inclusion command or drawing code, so no image is delivered and the package declares no asset dependency.  Editorial formatting only, in addition: an AMS \texttt{amsart} preamble that restates the article's own declarations and macros used by the retained lines, namely the environments \texttt{definition}, \texttt{proposition} on separate counters, together with the standard packages those lines need.  The retained environments print in this document as Definition 1, Proposition 1 in order of appearance, whereas the article prints the same items with the numbering of its own full document.  The complete native source of the article as distributed in the arXiv e-print for this exact version is bundled unchanged as \texttt{sources/191956/source0.tex}.
\begin{definition} \label{def:reduced_form}
   For any morphism $f\in \VI(n,a+1)$, define the morphism $\overline{f}\in \VI(n,a+1)$ as follows:
   \begin{itemize}
      \item if $\beta_f=0$, then $\overline{f}=f$;
      \item if $\beta_f=(\beta_1 \ \beta_2 \ \ldots \ \beta_n)$ with $\beta_1=\cdots=\beta_{j-1}=0$ and $\beta_j\neq 0$, then $\overline{f}=\sigma(c)f$ where $c$ is the unique vector in $\F^a$ such that the $j$-th column of $\gamma_{\sigma(c)f}$ is $0$.
   \end{itemize}    
\end{definition}

\begin{proposition} \label{prop:barshift1}
   Let $n\in \N$. 
   
   If $n=0$, then 
   \[ \SSigma M^{\VI}(n) \cong M^{\VI}(n). \]    
   
   If $n\geqslant 1$, then 
   \[ \SSigma M^{\VI}(n) \cong M^{\VI}(n) \oplus \left( M^{\VI}(n-1)^{\oplus (q^n-1)} \right). \]
\end{proposition}

\begin{proof}
   The $n=0$ case is trivial so assume that $n\geqslant 1$. 

   Let $a\in \N$ with $a\geqslant n-1$. For each row vector $\beta: \F^n \to \F$, let $P(\beta)_a$ be the $\Bbbk$-submodule of $(\Sigma M^{\VI}(n))_a$ generated by all $f\in \VI(n,a+1)$ such that $\beta_f=\beta$ and $\overline{f}=f$ (see Definition \ref{def:reduced_form}). In other words:
   \begin{itemize}
      \item if $\beta=0$, then $P(\beta)_a$ is generated by all $f\in \VI(n,a+1)$ such that $\beta_f=0$;
      \item if $\beta=(\beta_1 \ \beta_2 \ \ldots \ \beta_n)$ with $\beta_1=\cdots=\beta_{j-1}=0$ and $\beta_j\neq 0$, then $P(\beta)_a$ is generated by all $f\in \VI(n,a+1)$ such that $\beta_f=\beta$ and the $j$-th column of $\gamma_f$ is $0$.
   \end{itemize}    
   This defines a $\VI$-submodule $P(\beta)$ of $\Sigma M^{\VI}(n)$. It is easy to see that:
   \begin{itemize}
      \item if $\beta=0$, then $P(\beta)$ is isomorphic to $M^{\VI}(n)$;
      \item if $\beta\neq 0$, then $P(\beta)$ is isomorphic to $M^{\VI}(n-1)$.
   \end{itemize} 
   Let $P$ be the sum of $P(\beta)$ over all row vectors $\beta: \F^n \to \F$. Then 
   \[ P \cong M^{\VI}(n) \oplus \left( M^{\VI}(n-1)^{\oplus (q^n-1)} \right). \] 
   
   Let $\eta$ be the composition of the inclusion map $P\to \Sigma M(n)$ and the quotient map $\Sigma M(n) \to \SSigma M(n)$. It suffices to show that $\eta$ is an isomorphism. Let $a\in \N$ with $a\geqslant n-1$. We need to show that $\eta_a : P_a \to (\SSigma M(n))_a$ is bijective. 

   Let $\zeta_a: (\Sigma M(n))_a \to P_a$ be the $\Bbbk$-linear map sending each morphism $f\in \VI(n,a+1)$ to $\overline{f}$. Then $\zeta_a$ factors through $(\SSigma M(n))_a$ to give a $\Bbbk$-linear map $\overline{\zeta}_a: (\SSigma M(n))_a \to P_a$. It is plain that $\eta_a$ and $\overline{\zeta}_a$ are inverses of each other, hence $\eta_a$ is bijective.
\end{proof}

\end{document}
